\documentclass[../../main.tex]{subfiles}
\graphicspath{{\subfix{../../image/}}}

\begin{document}

\section{习题}\label{section:ch02-exercises}

在以下习题中，除非另作说明，所有给定函数均假定光滑。

\begin{exercise}\label{exercise:ch02-1}
写出一个函数 $u$ 求解初值问题
\begin{align*}
\begin{cases}
u_t+b\cdot Du+cu=0 & \text{在 }\mathbb{R}^n\times(0,\infty)\text{ 中},\\
u=g & \text{在 }\mathbb{R}^n\times\{t=0\}\text{ 上}
\end{cases}
\end{align*}
的显式公式。这里 $c\in\mathbb{R}$，$b\in\mathbb{R}^n$ 为常数。
\end{exercise}

\begin{exercise}\label{exercise:ch02-2}
证明Laplace方程 $\Delta u=0$ 在旋转下不变；即若 $O$ 为正交 $n\times n$ 矩阵，定义 $v(x)\triangleq u(Ox)$（$x\in\mathbb{R}^n$），则 $\Delta v=0$。
\end{exercise}

\begin{exercise}\label{exercise:ch02-3}
修改平均值公式的证明，证明对 $n\geqslant 3$，
\begin{align*}
u(0)=\fint_{\partial B(0,r)}g\,\mathrm{d}S+\frac{1}{n(n-2)\alpha(n)}\int_{B(0,r)}\left(\frac{1}{|x|^{n-2}}-\frac{1}{r^{n-2}}\right)f\,\mathrm{d}x,
\end{align*}
只要
\begin{align*}
\begin{cases}
-\Delta u=f & \text{在 }B_0(0,r)\text{ 中},\\
u=g & \text{在 }\partial B(0,r)\text{ 上}.
\end{cases}
\end{align*}
\end{exercise}

\begin{exercise}\label{exercise:ch02-4}
直接证明：若 $u\in C^2(U)\cap C(\overline{U})$ 在有界开集 $U$ 内调和，则 $\max_{\overline{U}}u=\max_{\partial U}u$。（提示：对 $\varepsilon>0$ 定义 $u_\varepsilon\triangleq u+\varepsilon|x|^2$，并证明 $u_\varepsilon$ 不可能在 $\overline{U}$ 的内部点达到其最大值。）
\end{exercise}

\begin{exercise}\label{exercise:ch02-5}
我们称 $v\in C^2(\overline{U})$ 为次调和的，若在 $U$ 中 $-\Delta v\leqslant 0$。
\begin{enumerate}[(a)]
\item 证明：对次调和函数 $v$，有 $v(x)\leqslant\fint_{B(x,r)}v\,\mathrm{d}y$，对一切 $B(x,r)\subset U$。
\item 证明：因而若 $U$ 有界，则 $\max_{\overline{U}}v=\max_{\partial U}v$。
\item 设 $\phi:\mathbb{R}\to\mathbb{R}$ 光滑且凸。设 $u$ 调和，$v\triangleq\phi(u)$。证明 $v$ 次调和。
\item 证明：只要 $u$ 调和，$v\triangleq|Du|^2$ 次调和。
\end{enumerate}
\end{exercise}

\begin{exercise}\label{exercise:ch02-6}
设 $U$ 为 $\mathbb{R}^n$ 的有界开子集。证明存在常数 $C$（仅依赖于 $U$），使只要 $u$ 是
\begin{align*}
\begin{cases}
-\Delta u=f & \text{在 }U\text{ 中},\\
u=g & \text{在 }\partial U\text{ 上}
\end{cases}
\end{align*}
的光滑解，就有 $\max_{\overline{U}}|u|\leqslant C\left(\max_{\partial U}|g|+\max_{\overline{U}}|f|\right)$。（提示：对 $\lambda\triangleq\max_{\overline{U}}|f|$，$-\Delta\left(u+\dfrac{|x|^2}{2n}\lambda\right)\leqslant 0$。）
\end{exercise}

\begin{exercise}\label{exercise:ch02-7}
用球的 Poisson 公式证明：只要 $u$ 在 $B_0(0,r)$ 中正且调和，就有
\begin{align*}
\frac{r^{n-2}}{r-|x|}(r+|x|)^{n-1}u(0)\leqslant u(x)\leqslant\frac{r^{n-2}}{r+|x|}(r-|x|)^{n-1}u(0).
\end{align*}
这是 Harnack 不等式的显式形式。
\end{exercise}

\begin{exercise}\label{exercise:ch02-8}
证明 \S2.2.4 定理\cref{theorem:ch02-poisson-ball}。（提示：由于 $u\equiv 1$ 在 $g\equiv 1$ 时解\eqref{align:ch02-ball-r-bvp}，理论自动给出对每个 $x\in B_0(0,1)$，$\int_{\partial B(0,1)}K(x,y)\,\mathrm{d}S(y)=1$。）
\end{exercise}

\begin{exercise}\label{exercise:ch02-9}
设 $u$ 是由半空间 Poisson 公式给出的
\begin{align*}
\begin{cases}
\Delta u=0 & \text{在 }\mathbb{R}^n_+\text{ 中},\\
u=g & \text{在 }\partial\mathbb{R}^n_+\text{ 上}
\end{cases}
\end{align*}
的解。设 $g$ 有界，且当 $x\in\partial\mathbb{R}^n_+$、$|x|\leqslant 1$ 时 $g(x)=|x|$。证明 $Du$ 在 $x=0$ 附近无界。（提示：估计 $\dfrac{u(\lambda e_n)-u(0)}{\lambda}$。）
\end{exercise}

\begin{exercise}[反射原理]\label{exercise:ch02-10}
\begin{enumerate}[(a)]
\item 设 $U_+$ 表示开半球 $\{x\in\mathbb{R}^n\mid |x|<1,\ x_n>0\}$。设 $u\in C^2(U_+)$ 在 $U_+$ 中调和，且在 $\partial U_+\cap\{x_n=0\}$ 上 $u=0$。令
\begin{align*}
v(x)\triangleq\begin{cases}
u(x) & \text{若 }x_n\geqslant 0,\\
-u(x_1,\cdots,x_{n-1},-x_n) & \text{若 }x_n<0
\end{cases}
\end{align*}
对 $x\in U=B_0(0,1)$。证明 $v\in C^2(U)$，从而 $v$ 在 $U$ 内调和。
\item 现只设 $u\in C^2(U_+)\cap C(\overline{U_+})$。证明 $v$ 在 $U$ 内调和。（提示：用球的 Poisson 公式。）
\end{enumerate}
\end{exercise}

\begin{exercise}[Laplace方程的 Kelvin 变换]\label{exercise:ch02-11}函数 $u:\mathbb{R}^n\to\mathbb{R}$ 的 Kelvin 变换 $\mathcal{K}u=\overline{u}$ 为
\begin{align*}
\overline{u}(x)\triangleq u(\overline{x})|\overline{x}|^{n-2}=u(x/|x|^2)|x|^{2-n} \qquad (x\neq 0),
\end{align*}
其中 $\overline{x}=x/|x|^2$。证明：若 $u$ 调和，则 $\overline{u}$ 也调和。（提示：先证 $D_x\overline{x}(D_x\overline{x})^T=|\overline{x}|^4 I$。映射 $x\mapsto\overline{x}$ 是共形的，即保角的。）
\end{exercise}

\begin{exercise}\label{exercise:ch02-12}
设 $u$ 光滑，在 $\mathbb{R}^n\times(0,\infty)$ 中解 $u_t-\Delta u=0$。
\begin{enumerate}[(a)]
\item 证明 $u_\lambda(x,t)\triangleq u(\lambda x,\lambda^2 t)$ 对每个 $\lambda\in\mathbb{R}$ 也解热方程。
\item 用 (a) 证明 $v(x,t)\triangleq x\cdot Du(x,t)+2t\,u_t(x,t)$ 也解热方程。
\end{enumerate}
\end{exercise}

\begin{exercise}\label{exercise:ch02-13}
设 $n=1$，$u(x,t)=v\!\left(\dfrac{x}{\sqrt{t}}\right)$。
\begin{enumerate}[(a)]
\item 证明 $u_t=u_{xx}$ 当且仅当
\begin{align*}
v''+\frac{z}{2}v'=0. 
\end{align*}
证明 $(*)$ 的通解为 $v(z)=c\displaystyle\int_0^z e^{-s^2/4}\,\mathrm{d}s+d$。
\item 对 $u(x,t)=v\!\left(\dfrac{x}{\sqrt{t}}\right)$ 关于 $x$ 求导，并适当选取常数 $c$，得到 $n=1$ 的基本解 $\Phi$。解释为何这一程序产生基本解。（提示：$u$ 的初始条件是什么？）
\end{enumerate}
\end{exercise}

\begin{exercise}\label{exercise:ch02-14}
写出
\begin{align*}
\begin{cases}
u_t-\Delta u+cu=f & \text{在 }\mathbb{R}^n\times(0,\infty)\text{ 中},\\
u=g & \text{在 }\mathbb{R}^n\times\{t=0\}\text{ 上}
\end{cases}
\end{align*}
一个解的显式公式，其中 $c\in\mathbb{R}$。
\end{exercise}

\begin{exercise}\label{exercise:ch02-15}
给定 $g:[0,\infty)\to\mathbb{R}$，$g(0)=0$，推导公式
\begin{align*}
u(x,t)=\frac{x}{\sqrt{4\pi}}\int_0^t\frac{1}{(t-s)^{3/2}}e^{-x^2/(4(t-s))}g(s)\,\mathrm{d}s
\end{align*}
为初/边值问题
\begin{align*}
\begin{cases}
u_t-u_{xx}=0 & \text{在 }\mathbb{R}_+\times(0,\infty)\text{ 中},\\
u=0 & \text{在 }\mathbb{R}_+\times\{t=0\}\text{ 上},\\
u=g & \text{在 }\{x=0\}\times[0,\infty)\text{ 上}
\end{cases}
\end{align*}
的一个解。（提示：令 $v(x,t)\triangleq u(x,t)-g(t)$，并把 $v$ 奇延拓到 $\{x<0\}$。）
\end{exercise}

\begin{exercise}\label{exercise:ch02-16}
直接证明：若 $U$ 有界，$u\in C^2_1(U_T)\cap C(\overline{U_T})$ 解热方程，则 $\max_{\overline{U_T}}u=\max_{\Gamma_T}u$。（提示：对 $\varepsilon>0$ 定义 $u_\varepsilon\triangleq u-\varepsilon t$，并证明 $u_\varepsilon$ 不可能在 $U_T$ 内一点达到它在 $\overline{U_T}$ 上的最大值。）
\end{exercise}

\begin{exercise}\label{exercise:ch02-17}
我们称 $v\in C^2_1(U_T)$ 为热方程的\textbf{次解}，若在 $U_T$ 中 $v_t-\Delta v\leqslant 0$。
\begin{enumerate}[(a)]
\item 证明：对次解 $v$，有 $v(x,t)\leqslant\dfrac{1}{4r^n}\displaystyle\iint_{E(x,t;r)}v(y,s)\dfrac{|x-y|^2}{(t-s)^2}\,\mathrm{d}y\,\mathrm{d}s$，对一切 $E(x,t;r)\subset U_T$。
\item 证明：因而 $\max_{\overline{U_T}}v=\max_{\Gamma_T}v$。
\item 设 $\phi:\mathbb{R}\to\mathbb{R}$ 光滑且凸。设 $u$ 解热方程，$v\triangleq\phi(u)$。证明 $v$ 是次解。
\item 证明：只要 $u$ 解热方程，$v\triangleq|Du|^2+u_t^2$ 是次解。
\end{enumerate}
\end{exercise}

\begin{exercise}\label{exercise:ch02-18}
设 $u$ 解初值问题
\begin{align*}
\begin{cases}
u_{tt}-\Delta u=0 & \text{在 }\mathbb{R}^n\times(0,\infty)\text{ 中},\\
u=0,\ u_t=h & \text{在 }\mathbb{R}^n\times\{t=0\}\text{ 上}.
\end{cases}
\end{align*}
证明 $v\triangleq u_t$ 解
\begin{align*}
\begin{cases}
v_{tt}-\Delta v=0 & \text{在 }\mathbb{R}^n\times(0,\infty)\text{ 中},\\
v=h,\ v_t=0 & \text{在 }\mathbb{R}^n\times\{t=0\}\text{ 上}.
\end{cases}
\end{align*}
这是 Stokes 法则。
\end{exercise}

\begin{exercise}\label{exercise:ch02-19}
\begin{enumerate}[(a)]
\item 证明 PDE $u_{xy}=0$ 的通解为 $u(x,y)=F(x)+G(y)$，其中 $F,G$ 为任意函数。
\item 作变量代换 $\xi=x+t$，$\eta=x-t$，证明 $u_{tt}-u_{xx}=0$ 当且仅当 $u_{\xi\eta}=0$。
\item 用 (a)、(b) 重新导出 d'Alembert 公式。
\item 初值 $g,h$ 满足什么条件时，解 $u$ 是右行波？左行波？
\end{enumerate}
\end{exercise}

\begin{exercise}\label{exercise:ch02-20}
设对某个衰减函数 $\alpha=\alpha(r)$ 与延迟函数 $\beta=\beta(r)\geqslant 0$，对一切波形 $\phi$，波动方程在 $(\mathbb{R}^n-\{0\})\times\mathbb{R}$ 中存在形如
\begin{align*}
u(x,t)=\alpha(r)\phi(t-\beta(r))
\end{align*}
的解，其中 $r=|x|$，并设 $\beta(0)=0$。证明这只有在 $n=1$ 或 $n=3$ 时才可能，并计算函数 $\alpha,\beta$ 的形式。（T. Morley, SIAM Review 27 (1985), 69--71）
\end{exercise}

\begin{exercise}\label{exercise:ch02-21}
\begin{enumerate}[(a)]
\item 设 $\mathbf{E}=(E_1,E_2,E_3)$，$\mathbf{B}=(B_1,B_2,B_3)$ 解 Maxwell 方程
\begin{align*}
\mathbf{E}_t=\operatorname{curl}\mathbf{B},\qquad \mathbf{B}_t=-\operatorname{curl}\mathbf{E},\qquad \operatorname{div}\mathbf{B}=\operatorname{div}\mathbf{E}=0.
\end{align*}
证明 $\mathbf{E}_{tt}-\Delta\mathbf{E}=0$，$\mathbf{B}_{tt}-\Delta\mathbf{B}=0$。
\item 设 $\mathbf{u}=(u_1,u_2,u_3)$ 解线性弹性力学的发展方程 $\mathbf{u}_{tt}-\mu\Delta\mathbf{u}-(\lambda+\mu)D(\operatorname{div}\mathbf{u})=0$（在 $\mathbb{R}^3\times(0,\infty)$ 中）。证明 $w\triangleq\operatorname{div}\mathbf{u}$ 与 $\mathbf{w}\triangleq\operatorname{curl}\mathbf{u}$ 各自解波动方程，但传播速度不同。
\end{enumerate}
\end{exercise}

\begin{exercise}\label{exercise:ch02-22}
设 $u$ 表示沿实线以速度 1 向右运动的粒子密度，$v$ 表示以速度 1 向左运动的粒子密度。若向右运动的粒子以速率 $d>0$ 随机变为向左运动，反之亦然，则有 PDE 组
\begin{align*}
\begin{cases}
u_t+u_x=d(v-u),\\
v_t-v_x=d(u-v).
\end{cases}
\end{align*}
证明 $w\triangleq u$ 与 $w\triangleq v$ 都解电报方程 $w_{tt}+2d w_t-w_{xx}=0$。
\end{exercise}

\begin{exercise}\label{exercise:ch02-23}
设 $S$ 表示 $\mathbb{R}\times(0,\infty)$ 中以 $(0,1),(1,2),(0,3),(-1,2)$ 为顶点的正方形。定义
\begin{align*}
f(x,t)\triangleq
\begin{cases}
-1 & \text{若 }(x,t)\in S\cap\{t>x+2\},\\
1 & \text{若 }(x,t)\in S\cap\{t<x+2\},\\
0 & \text{其他}.
\end{cases}
\end{align*}
设 $u$ 解
\begin{align*}
\begin{cases}
u_{tt}-u_{xx}=f & \text{在 }\mathbb{R}\times(0,\infty)\text{ 中},\\
u=0,\ u_t=0 & \text{在 }\mathbb{R}\times\{t=0\}\text{ 上}.
\end{cases}
\end{align*}
描述时刻 $t>3$ 时 $u$ 的形状。（J. G. Kingston, SIAM Review 30 (1988), 645--649）
\end{exercise}

\begin{exercise}[能量均分]\label{exercise:ch02-24}设 $u$ 解一维波动方程的初值问题：
\begin{align*}
\begin{cases}
u_{tt}-u_{xx}=0 & \text{在 }\mathbb{R}\times(0,\infty)\text{ 中},\\
u=g,\ u_t=h & \text{在 }\mathbb{R}\times\{t=0\}\text{ 上}.
\end{cases}
\end{align*}
设 $g,h$ 具有紧支集。动能为 $k(t)\triangleq\dfrac{1}{2}\displaystyle\int_{-\infty}^{\infty}u_t^2(x,t)\,\mathrm{d}x$，势能为 $p(t)\triangleq\dfrac{1}{2}\displaystyle\int_{-\infty}^{\infty}u_x^2(x,t)\,\mathrm{d}x$。证明
\begin{enumerate}[(a)]
\item $k(t)+p(t)$ 关于 $t$ 为常数；
\item 对一切足够大的时刻 $t$，$k(t)=p(t)$。
\end{enumerate}
\end{exercise}

\end{document}